\documentclass{ctexart}
% claim-forward 模块中文 fixture（latex-thesis-zh）。
% Case A–E 为五个观察码各一正例；Case F–I 为必须保持沉默的边界。
% 供 evals/evals.json id 49 与 tests/skills/latex_thesis_zh/test_claim_forward_zh.py 使用。
\begin{document}

\chapter{摘要}
本文提出一种轻量级跟踪方法，在三个室内基准上将准确率提高12\%。
该方法在单块 GPU 上达到每秒210帧。

\chapter{绪论}
% Case A：CF-DISCLAIM——段首以“本文不试图”开头，主张在后。
本文不试图解决通用场景理解问题。
本文提出一种轻量级跟踪方法，在室内基准上将准确率提高12\%。

\section{主要贡献}
% Case C：CF-CAVEAT-POS——关于自身工作的限制句写在其限定的主张句之前。
尽管本文的评估范围限于室内场景，相关结论应在该范围内理解。
本文提出的跟踪方法相比已发表的最优方法将延迟降低40\%。

\chapter{方法设计}
% Case D：CF-HEDGE-STACK——同一主张句上叠加三个 hedge。
本文提出的方法可能在一定程度上或许能够提升鲁棒性。
匹配模块采用固定的八帧窗口。

\chapter{实验结果}
% Case B：CF-SELFWEAK——对自身结果使用自我削弱搭配。
本文方法在标准划分上达到94.1\%的准确率。
遗憾的是，本文方法在长尾划分上仍明显落后于理想上界。

% Case H：边界——裸“仅为”与“不仅”不是自我削弱。
本文方法的误差仅为0.018，不仅优于基线，而且内存占用更低。

% Case I：边界——“尚未”是问题陈述用语（摘要痛点词），不报。
现有方法尚未解决长序列漂移问题。

\chapter{讨论}
% Case F：边界——引用他人工作时可以说其不足。
已有方法\cite{tracker2019}未能在普通硬件上达到实时速率。
该类方法在长序列上仍明显落后于离线方法。

\section{本文的不足}
% Case G：边界——不足小节本来就写限制。
本文的评估未能覆盖强光照变化的室外场景。
本文仅在三个室内基准上给出结果。

\chapter{结论与展望}
本文提出了一种轻量级跟踪方法，并在三个室内基准上验证了12\%的准确率提升。
该设计将单帧开销控制在5毫秒以内。

% Case E：CF-CLOSE-NEG——末段以负面判定收尾，无展望方向。
本文提出的框架在单块 GPU 上达到实时速率。
然而，该框架无法处理帧率可变的流式输入。

\end{document}
